\section{Rotary Encoders}

\subsection{Overview}

Two rotary encoders are used in this design. One controls the alert buzzer volume and
the other controls the display brightness. Either encoder may be temporarily
reassigned by the application script to set other parameters as needed. The encoders
do not include push-button switches.

A rotary encoder is a digital position sensor. Unlike a potentiometer, it has no
end stops and produces no analog output. Instead, it produces a series of digital
pulses as the shaft is rotated, with the direction of rotation encoded in the phase
relationship between two output signals. Because the encoder has no absolute position,
the application maintains a software counter that is incremented or decremented on
each detent step.

\subsection{How a Rotary Encoder Works}

Inside the encoder are two switches, conventionally labelled A and B (sometimes CLK
and DT). The switches are arranged so that they open and close in a specific sequence
depending on the direction of rotation. As the shaft turns clockwise, A transitions
before B. As the shaft turns counter-clockwise, B transitions before A. The
relationship is a two-bit Gray code sequence:

\begin{table}[h]
\centering
\begin{tabular}{ccc}
\hline
\textbf{Step} & \textbf{A} & \textbf{B} \\
\hline
0 & 1 & 1 \\
1 & 0 & 1 \\
2 & 0 & 0 \\
3 & 1 & 0 \\
\hline
\end{tabular}
\caption{Gray code sequence for clockwise rotation. Counter-clockwise reverses the order.}
\label{tab:encoder_gray}
\end{table}

The simplest detection strategy is to trigger an interrupt on the falling edge of
signal A, then immediately read the state of signal B. If B is high at that moment,
rotation is clockwise; if B is low, rotation is counter-clockwise. This single-edge
approach is sufficient for a hand-turned control knob where speed is low.

\subsection{Contact Bounce and Debouncing}

The switches inside a mechanical rotary encoder bounce as they make and break contact,
producing multiple spurious transitions per detent step. Without debouncing, the
interrupt handler fires several times per click and the counter accumulates incorrect
values.

Two debouncing strategies are available:

\begin{itemize}
    \item \textbf{Hardware debouncing:} A small capacitor (typically \SI{10}{\nano\farad}
    to \SI{100}{\nano\farad}) from each encoder pin to ground, combined with a pull-up
    resistor (\SI{10}{\kilo\ohm}), forms an RC low-pass filter that smooths the bounce
    transients. This is the preferred approach for a production design as it requires no
    timing logic in software.

    \item \textbf{Software debouncing:} Record the timestamp of each interrupt and
    ignore subsequent triggers that arrive within a debounce window (typically
    \SI{5}{\milli\second} to \SI{50}{\milli\second}). The RPi.GPIO
    \texttt{add\_event\_detect()} call includes a built-in \texttt{bouncetime}
    parameter for exactly this purpose.
\end{itemize}

Both approaches may be combined. Hardware debouncing is recommended with a
\SI{100}{\nano\farad} ceramic capacitor on each encoder pin.

\subsection{Wiring}

Each encoder has two signal pins (A and B), a common pin, and optionally a supply
pin. The common is connected to ground. Each signal pin is pulled up to \SI{3.3}{\volt}
through a \SI{10}{\kilo\ohm} resistor (the Raspberry Pi internal pull-up may be used
instead, activated in software). A \SI{100}{\nano\farad} ceramic capacitor from each
signal pin to ground provides hardware debouncing.

\begin{verbatim}
  3.3V ---[ 10k ]---+--- Encoder A pin
                    |
                  100nF
                    |
                   GND
                   
  3.3V ---[ 10k ]---+--- Encoder B pin
                    |
                  100nF
                    |
                   GND

  Encoder Common --- GND
\end{verbatim}

\subsection{Pin Restrictions on the Raspberry Pi Zero}

All GPIO pins on the Raspberry Pi Zero support interrupts via the RPi.GPIO library.
There is no hardware restriction limiting interrupts to specific pins, unlike some
microcontrollers. The Linux kernel \texttt{gpio} subsystem services all GPIO interrupts
through a unified interface.

Practical limits to be aware of:

\begin{itemize}
    \item \textbf{Timing jitter:} Because the Raspberry Pi runs Linux, interrupt
    latency is not deterministic. Typical latency is \SI{1}{\milli\second} to
    \SI{5}{\milli\second} under light load, but may spike under heavy CPU or I/O
    activity. For a hand-turned knob this is entirely acceptable; for high-speed
    pulse counting it would not be.

    \item \textbf{Number of interrupts:} There is no fixed limit on the number of
    simultaneous GPIO interrupts. Each \texttt{add\_event\_detect()} call starts a
    kernel-level watch on that pin. In practice, a handful of interrupts is well
    within the capability of the system.

    \item \textbf{Power:} GPIO pins source or sink a maximum of \SI{16}{\milli\ampere}
    individually, with a total bank limit. Encoder signal pins draw negligible current
    when used as inputs with pull-ups; no power budget concern arises.

    \item \textbf{Voltage:} All GPIO inputs must remain within \SI{0}{\volt} to
    \SI{3.3}{\volt}. The encoders are passive switches to ground and are inherently
    safe provided the pull-ups are connected to \SI{3.3}{\volt} and not \SI{5}{\volt}.
\end{itemize}

\subsection{Setting Up Interrupts in Python}

The RPi.GPIO library provides \texttt{add\_event\_detect()} to register a callback
function that is invoked when a specified edge transition is detected on a GPIO pin.
The callback runs in a separate thread managed by the library, but in a simple
synchronous script the main loop can simply read shared variables that the callback
updates.

The general pattern for a single interrupt is:

\begin{verbatim}
import RPi.GPIO as GPIO

GPIO.setmode(GPIO.BCM)

PIN = 17
GPIO.setup(PIN, GPIO.IN, pull_up_down=GPIO.PUD_UP)

def on_edge(channel):
    # called automatically when the edge is detected
    pass

GPIO.add_event_detect(
    PIN,
    GPIO.FALLING,
    callback=on_edge,
    bouncetime=10        # ignore re-triggers within 10 ms
)
\end{verbatim}

\subsection{Two Encoders with Multiple Interrupts}

For two encoders, four GPIO pins are used (A and B for each encoder). An interrupt
is registered on the A pin of each encoder. When the interrupt fires, the handler
reads the B pin to determine direction and updates the appropriate software counter.

\begin{verbatim}
import RPi.GPIO as GPIO
import time

# --- Pin assignments ---
ENC1_A = 17    # Volume encoder
ENC1_B = 27
ENC2_A = 22    # Brightness encoder
ENC2_B = 23

# --- Shared state ---
volume     = 50    # 0-100
brightness = 50    # 0-100

GPIO.setmode(GPIO.BCM)

for pin in [ENC1_A, ENC1_B, ENC2_A, ENC2_B]:
    GPIO.setup(pin, GPIO.IN, pull_up_down=GPIO.PUD_UP)


def on_volume(channel):
    global volume
    direction = 1 if GPIO.input(ENC1_B) else -1
    volume = max(0, min(100, volume + direction))


def on_brightness(channel):
    global brightness
    direction = 1 if GPIO.input(ENC2_B) else -1
    brightness = max(0, min(100, brightness + direction))


GPIO.add_event_detect(ENC1_A, GPIO.FALLING,
                      callback=on_volume,     bouncetime=10)
GPIO.add_event_detect(ENC2_A, GPIO.FALLING,
                      callback=on_brightness, bouncetime=10)


def main():
    print("Turning encoders. Ctrl-C to exit.")
    try:
        while True:
            # Main loop runs normally; callbacks fire asynchronously
            print(f"Volume: {volume:3d}   Brightness: {brightness:3d}",
                  end="\r")
            time.sleep(0.1)
    except KeyboardInterrupt:
        pass
    finally:
        GPIO.cleanup()


if __name__ == "__main__":
    main()
\end{verbatim}

\noindent Key points in this implementation:

\begin{itemize}
    \item The B pin of each encoder is read inside the interrupt callback, not
    registered as a separate interrupt. Only the A pin triggers the interrupt.
    \item The \texttt{bouncetime=10} argument suppresses re-triggers within
    \SI{10}{\milli\second}, providing software debouncing in addition to any
    hardware RC filtering.
    \item The \texttt{global} keyword allows the callback (running in the RPi.GPIO
    interrupt thread) to modify the counter variables read by the main loop.
    \item The counters are clamped to their valid ranges inside each callback.
    \item \texttt{GPIO.cleanup()} is called on exit to release all pin reservations
    and remove event detection.
\end{itemize}

\subsection{Temporarily Reassigning an Encoder}

Because the encoder counters are simple integer variables, reassigning an encoder
to a different parameter requires only redirecting which variable the callback
modifies. The cleanest approach is to use a function that unregisters the existing
callback and registers a new one:

\begin{verbatim}
def reassign_encoder1(new_variable: list, label: str) -> None:
    """
    Temporarily redirect encoder 1 to control a different parameter.
    new_variable must be a mutable container, e.g. a list of one int.
    """
    GPIO.remove_event_detect(ENC1_A)

    def on_reassigned(channel):
        direction = 1 if GPIO.input(ENC1_B) else -1
        new_variable[0] = max(0, min(100, new_variable[0] + direction))

    GPIO.add_event_detect(ENC1_A, GPIO.FALLING,
                          callback=on_reassigned, bouncetime=10)
    print(f"Encoder 1 now controls: {label}")
\end{verbatim}

\noindent To restore the original volume function, call \texttt{reassign\_encoder1}
again with the \texttt{volume} container and label.

\subsection{Adafruit 1651 Touchscreen Interrupt (TSC2007)}

The Adafruit 1651 display uses the \textbf{TSC2007} resistive touch controller,
which communicates over I2C. The TSC2007 provides a \textbf{PENIRQ} pin (active
low) that is asserted whenever the screen is touched. Connecting this pin to a
GPIO and registering a falling-edge interrupt allows the application to respond to
touch events immediately, without polling the I2C bus in a continuous loop.

\begin{verbatim}
TOUCH_IRQ = 24    # PENIRQ from TSC2007

GPIO.setup(TOUCH_IRQ, GPIO.IN, pull_up_down=GPIO.PUD_UP)

def on_touch(channel):
    """Read touch coordinates from TSC2007 over I2C."""
    # Issue a read command to the TSC2007 over I2C here
    # e.g. using smbus2 or adafruit_tsc2007 library
    pass

GPIO.add_event_detect(TOUCH_IRQ, GPIO.FALLING,
                      callback=on_touch, bouncetime=50)
\end{verbatim}

\noindent The PENIRQ pin requires a pull-up to \SI{3.3}{\volt}; the internal
RPi pull-up activated by \texttt{GPIO.PUD\_UP} is sufficient. The PENIRQ line
is open-drain on the TSC2007 and safe to connect directly to the GPIO pin
provided the pull-up is to \SI{3.3}{\volt}.

\subsection{Parts List}

\begin{table}[h]
\centering
\begin{tabular}{llll}
\hline
\textbf{Reference} & \textbf{Description} & \textbf{Value / Part} & \textbf{Package} \\
\hline
ENC1, ENC2  & Rotary encoder, mechanical, no push button & e.g. Alps EC11  & Through-hole \\
R1--R4      & Signal pull-up resistors (A and B, each encoder) & \SI{10}{\kilo\ohm} & 0603 \\
C1--C4      & Hardware debounce capacitors (A and B, each encoder) & \SI{100}{\nano\farad} ceramic & 0603 \\
\hline
\end{tabular}
\caption{Rotary encoder components}
\label{tab:encoder_parts}
\end{table}
